\documentclass[../../main.tex]{subfiles}
\begin{document}

\section{带边流形}\label{section:ch01-带边流形}

在流形的许多重要应用中，最引人注目的是那些涉及积分的应用，我们会遇到一些本来会是光滑流形、却带有某种"边界"的空间。这类空间的简单例子包括 $\mathbb{R}$ 中的闭区间、$\mathbb{R}^{n}$ 中的闭球以及 $S^{n}$ 中的闭半球。为容纳这类空间，我们需要推广流形的定义。

这些空间中的点将有局部建模在 $\mathbb{R}^{n}$ 的开子集上、或建模在闭 $n$ 维上半空间 $H^{n}\subset\mathbb{R}^{n}$ 的开子集上的邻域，$H^{n}$ 定义为
\begin{align*}
H^{n}=\{(x^{1},\dots,x^{n})\in\mathbb{R}^{n}:x^{n}\geqslant 0\}.
\end{align*}
我们用记号 $\operatorname{Int}H^{n}$ 和 $\partial H^{n}$ 分别表示 $H^{n}$ 作为 $\mathbb{R}^{n}$ 子集的内部和边界。当 $n>0$ 时，这意味着
\begin{align*}
\operatorname{Int}H^{n}=\{(x^{1},\dots,x^{n})\in\mathbb{R}^{n}:x^{n}>0\},\qquad
\partial H^{n}=\{(x^{1},\dots,x^{n})\in\mathbb{R}^{n}:x^{n}=0\}.
\end{align*}
在 $n=0$ 的情形，$H^{0}=\mathbb{R}^{0}=\{0\}$，所以 $\operatorname{Int}H^{0}=\mathbb{R}^{0}$，$\partial H^{0}=\varnothing$。

一个 $n$ 维拓扑带边流形是一个第二可数的 Hausdorff 空间 $M$，其中每个点都有一个邻域同胚于 $\mathbb{R}^{n}$ 的某个开子集、或同胚于 $H^{n}$ 的某个（相对）开子集（图~\cref{figure:ch01-1.11}）。一个开子集 $U\subset M$ 连同一个到 $\mathbb{R}^{n}$ 或 $H^{n}$ 的开子集为同胚的映射 $\varphi:U\to\mathbb{R}^{n}$，将像流形的情形那样称为 $M$ 的一个图。当需要区分时，若 $\varphi(U)$ 是 $\mathbb{R}^{n}$ 的一个开子集（包括 $H^{n}$ 中不与 $\partial H^{n}$ 相交的开子集的情形），我们称 $(U,\varphi)$ 为内图；若 $\varphi(U)$ 是 $H^{n}$ 的一个开子集且 $\varphi(U)\cap\partial H^{n}\neq\varnothing$，则称其为边界图。一个像为形如 $B_{r}(x)\cap H^{n}$（其中 $x\in\partial H^{n}$，$r>0$）的集合的边界图称为坐标半球。

\begin{figure}[ht]
\centering
\includegraphics[width=0.8\textwidth]{fig-01-11.png}
\caption{带边流形}\label{figure:ch01-1.11}
\end{figure}

称点 $p\in M$ 为 $M$ 的内点，如果它在某个内图的定义域中。称它为 $M$ 的边界点，如果它在某个把 $p$ 映到 $\partial H^{n}$ 的边界图的定义域中。$M$ 的边界（其所有边界点所成集合）记作 $\partial M$；类似地，它的内部（所有内点所成集合）记作 $\operatorname{Int}M$。

由定义可得，每个点 $p\in M$ 要么是内点、要么是边界点：若 $p$ 不是边界点，则它要么在某个内图的定义域中，要么在某个边界图 $(U,\varphi)$ 的定义域中而 $\varphi(p)\notin\partial H^{n}$，在后一情形 $\varphi$ 在 $U\cap\varphi^{-1}(\operatorname{Int}H^{n})$ 上的限制就是一个定义域包含 $p$ 的内图。然而，一个给定的点不可能同时关于一个图是内点、关于另一个图是边界点，这一点并不显然。事实上这不会发生，但证明需要比我们目前所掌握的更多工具。为方便起见，我们在此陈述该定理。

\begin{theorem}[边界的拓扑不变性]\label{theorem:ch01-边界拓扑不变性}
若 $M$ 是一个拓扑带边流形，则 $M$ 的每个点要么是边界点、要么是内点，但不同时是两者。因此 $\partial M$ 与 $\operatorname{Int}M$ 是两个不相交集合，其并为 $M$。
\end{theorem}

证明见习题~17-9。在本章稍后，我们将对光滑带边流形证明这一结果的一个较弱版本（定理~\cref{theorem:ch01-边界光滑不变性}），它对我们的大多数目的已经足够。

请务必注意区分术语"边界"和"内部"的这些新定义，与它们用来指拓扑空间某子集的边界和内部时的用法。一个带边流形在这个新意义下可能有非空边界，这与它作为某个其他拓扑空间的子集是否有边界无关。如果需要强调两种边界概念之间的区别，我们将酌情使用术语拓扑边界和流形边界。例如，闭单位球 $\overline{B^{n}}$ 是一个带边流形（见习题~1-11），其流形边界是 $S^{n-1}$。它作为 $\mathbb{R}^{n}$ 子集的拓扑边界恰好也是球面。然而，如果我们把 $\overline{B^{n}}$ 本身看作一个拓扑空间，则作为它自身的子集，它的拓扑边界是空集。而如果把它看作 $\mathbb{R}^{n+1}$ 的子集（以显然方式把 $\mathbb{R}^{n}$ 看作 $\mathbb{R}^{n+1}$ 的子集），它的拓扑边界就是整个 $\overline{B^{n}}$。注意 $H^{n}$ 本身是一个带边流形，其流形边界与它作为 $\mathbb{R}^{n}$ 子集的拓扑边界相同。$\mathbb{R}$ 中的每个区间都是一个 $1$ 维带边流形，其流形边界由它的端点（如果有）组成。

带边流形的命名是传统且确立已久的，但必须谨慎使用。尽管有其名称，带边流形一般并不是流形，因为边界点没有局部欧几里得邻域。（这是边界不变性定理的一个推论。）此外，带边流形的边界可能是空集——定义中没有任何要求边界非空。另一方面，一个流形也是一个边界为空的带边流形。因此，每个流形都是一个带边流形，而一个带边流形是流形当且仅当其边界为空（见下面的命题~\cref{proposition:ch01-带边流形基本性质}）。

尽管"带边流形"这一术语也涵盖流形，我们仍会常用诸如无边流形这样的冗余短语，如果我们想强调谈论的是原初意义下的流形；用带边或无边流形来指一个带边流形，如果我们想强调边界可能为空。（后一短语常在我们主要关心流形、但所讨论结果在更一般的带边流形情形下同样容易陈述和证明时出现。）注意，不加进一步限定的"流形"一词总指无边流形。在文献中，你还会遇到闭流形一词指无边紧流形，开流形一词指无边非紧流形。

\begin{proposition}\label{proposition:ch01-带边流形基本性质}
设 $M$ 是一个拓扑 $n$ 维带边流形。
\begin{enumerate}[(a)]
\item $\operatorname{Int}M$ 是 $M$ 的一个开子集，且是一个拓扑 $n$ 维无边流形。
\item $\partial M$ 是 $M$ 的一个闭子集，且是一个拓扑 $(n-1)$ 维无边流形。
\item $M$ 是一个拓扑流形当且仅当 $\partial M=\varnothing$。
\item 若 $n=0$，则 $\partial M=\varnothing$，且 $M$ 是一个 $0$ 维流形。
\end{enumerate}
\end{proposition}

\begin{exercise}\label{exercise:ch01-1.39}
证明上面的命题。在这个证明中，必要时可以使用边界拓扑不变性定理。哪些部分需要用到它？
\end{exercise}

我们在本章前面证明的流形拓扑性质，对带边流形有自然的推广，证明本质上与流形情形相同。为记录在案，我们在此陈述它们。

\begin{proposition}\label{proposition:ch01-带边流形拓扑性质}
设 $M$ 是一个拓扑带边流形。
\begin{enumerate}[(a)]
\item $M$ 有一个由预紧坐标球和坐标半球组成的可数基。
\item $M$ 是局部紧的。
\item $M$ 是仿紧的。
\item $M$ 是局部道路连通的。
\item $M$ 有可数多个分支，每个分支都是 $M$ 的一个开子集，且是一个连通的拓扑带边流形。
\item $M$ 的基本群是可数的。
\end{enumerate}
\end{proposition}

\begin{exercise}\label{exercise:ch01-1.41}
证明上面的命题。
\end{exercise}

\subsection{带边流形上的光滑结构}\label{subsection:ch01-带边流形光滑结构}

为看出如何在带边流形上定义光滑结构，回忆：从 $\mathbb{R}^{n}$ 的任意子集 $A$ 到 $\mathbb{R}^{k}$ 的映射称为光滑的，如果在 $A$ 的每一点的一个邻域内它都容许一个延拓，成为定义在 $\mathbb{R}^{n}$ 的某个开子集上的光滑映射（见附录~C，第~645 页）。因此，若 $U$ 是 $H^{n}$ 的一个开子集，映射 $F:U\to\mathbb{R}^{k}$ 光滑，是指对每个 $x\in U$，存在一个包含 $x$ 的开子集 $\widetilde{U}\subset\mathbb{R}^{n}$ 和一个光滑映射 $\widetilde{F}:\widetilde{U}\to\mathbb{R}^{k}$，使它在 $\widetilde{U}\cap H^{n}$ 上与 $F$ 一致（图~\cref{figure:ch01-1.12}）。若 $F$ 是这样一个映射，则 $F$ 在 $U\cap\operatorname{Int}H^{n}$ 上的限制在通常意义下是光滑的。由连续性，$F$ 在 $U\cap\partial H^{n}$ 各点处的所有偏导数都由它们在 $\operatorname{Int}H^{n}$ 中的值确定，因而特别地与延拓的选取无关。有一个事实（我们既不证明也不使用）：$F:U\to\mathbb{R}^{k}$ 在这个意义下光滑，当且仅当 $F$ 连续、$F|_{U\cap\operatorname{Int}H^{n}}$ 光滑，且 $F|_{U\cap\operatorname{Int}H^{n}}$ 的各阶偏导数都连续地延拓到整个 $U$。（一个方向是显然的；另一个方向依赖于 Émile Borel 的一个引理，该引理表明存在一个定义在下半空间的光滑函数，其在 $U\cap\partial H^{n}$ 上的导数都与 $F$ 的导数相匹配。参见，例如，[Hör90, 定理~1.2.6]。）

\begin{figure}[ht]
\centering
\includegraphics[width=0.8\textwidth]{fig-01-12.png}
\caption{$\mathbb{H}^{n}$ 开子集上映射的光滑性}\label{figure:ch01-1.12}
\end{figure}

例如，令 $B^{2}\subset\mathbb{R}^{2}$ 为开单位圆盘，$U=B^{2}\cap H^{2}$，定义 $f:U\to\mathbb{R}$ 为 $f(x,y)=\sqrt{1-x^{2}-y^{2}}$。因为 $f$（按同一公式）光滑延拓到整个 $B^{2}$，$f$ 是 $U$ 上的光滑函数。另一方面，虽然 $g(x,y)=\sqrt{y}$ 在 $U$ 上连续且在 $U\cap\operatorname{Int}H^{2}$ 中光滑，但它在 $\mathbb{R}^{2}$ 中原点的任意邻域内没有光滑延拓，因为当 $y\to 0$ 时 $\partial g/\partial y\to\infty$。因此 $g$ 在 $U$ 上不光滑。

现在设 $M$ 是一个拓扑带边流形。与流形情形一样，$M$ 的一个光滑结构定义为一个极大光滑图册——一族图，其定义域覆盖 $M$，且其转移映射（及其逆）在刚才描述的意义下光滑。赋予这样一个结构后，$M$ 称为一个光滑带边流形。每个光滑流形自动地是一个光滑带边流形（其边界为空）。

与光滑流形一样，若 $M$ 是一个光滑带边流形，给定光滑图册中的任意图称为 $M$ 的一个光滑图。$M$ 中的光滑坐标球、光滑坐标半球和正则坐标球按显然方式定义。此外，称子集 $B\subset M$ 为正则坐标半球，如果存在一个光滑坐标半球 $B'\supset\overline{B}$ 和一个光滑坐标映射 $\varphi:B'\to H^{n}$，使得对某对 $r'>r>0$ 有
\begin{align*}
\varphi(B)=B_{r}(0)\cap H^{n},\qquad
\varphi(\overline{B})=\overline{B_{r}(0)}\cap H^{n},\qquad
\varphi(B')=B_{r'}(0)\cap H^{n}.
\end{align*}

\begin{exercise}\label{exercise:ch01-1.42}
证明每个光滑带边流形都有一个由正则坐标球和半球组成的可数基。
\end{exercise}

\begin{exercise}\label{exercise:ch01-1.43}
证明光滑流形坐标图引理（引理~\cref{lemma:ch01-光滑流形坐标图引理}）在把"$\mathbb{R}^{n}$"替换为"$\mathbb{R}^{n}$ 或 $H^{n}$"、把"光滑流形"替换为"光滑带边流形"后仍然成立。
\end{exercise}

\begin{exercise}\label{exercise:ch01-1.44}
设 $M$ 是一个光滑 $n$ 维带边流形，$U$ 是 $M$ 的一个开子集。证明以下命题：
\begin{enumerate}[(a)]
\item $U$ 是一个拓扑 $n$ 维带边流形，而由 $M$ 的所有满足 $V\subset U$ 的光滑图 $(V,\varphi)$ 组成的图册在 $U$ 上定义了一个光滑结构。带这个拓扑和光滑结构，$U$ 称为一个开的带边子流形。
\item 若 $U\subset\operatorname{Int}M$，则 $U$ 实际上是一个（无边）光滑流形；在这种情形我们称它为 $M$ 的一个开子流形。
\item $\operatorname{Int}M$ 是 $M$ 的一个开子流形（无边）。
\end{enumerate}
\end{exercise}

关于光滑流形的一个重要结果不能直接推广到光滑带边流形，那就是有限乘积上光滑结构的构造（见例~\cref{example:ch01-乘积流形}）。因为半空间的乘积 $H^{n}\times H^{m}$ 本身不是半空间，光滑带边流形的有限乘积一般不能被看作光滑带边流形。（相反，它是光滑带角流形的一个例子，我们将在第~16 章研究。）然而，我们确实有下面的结果。

\begin{proposition}\label{proposition:ch01-乘积带边}
设 $M_{1},\dots,M_{k}$ 是光滑流形，$N$ 是一个光滑带边流形。则 $M_{1}\times\cdots\times M_{k}\times N$ 是一个光滑带边流形，且
\begin{align*}
\partial(M_{1}\times\cdots\times M_{k}\times N)=M_{1}\times\cdots\times M_{k}\times\partial N.
\end{align*}
\end{proposition}

\begin{proof}
习题~1-12。
\end{proof}

对光滑带边流形，下面的结果常常可以充分替代边界不变性定理。

\begin{theorem}[边界的光滑不变性]\label{theorem:ch01-边界光滑不变性}
设 $M$ 是一个光滑带边流形，$p\in M$。若存在 $M$ 的某个光滑图 $(U,\varphi)$ 使 $\varphi(U)\subset H^{n}$ 且 $\varphi(p)\in\partial H^{n}$，则每个定义域包含 $p$ 的光滑图也都如此。
\end{theorem}

\begin{proof}
假设相反，$p$ 既在一个光滑内图 $(U,\psi)$ 的定义域中，又在一个光滑边界图 $(V,\varphi)$ 的定义域中且 $\varphi(p)\in\partial H^{n}$。令 $\rho=\varphi\circ\psi^{-1}$ 表示转移映射；它是从 $\psi(U\cap V)$ 到 $\varphi(U\cap V)$ 的一个同胚。图的光滑相容性保证 $\rho$ 和 $\rho^{-1}$ 都光滑，其意义是：局部地，必要时它们可以延拓为定义在 $\mathbb{R}^{n}$ 开子集上的光滑映射。

写 $x_{0}=\psi(p)$，$y_{0}=\varphi(p)=\rho(x_{0})$。在 $\mathbb{R}^{n}$ 中存在 $y_{0}$ 的某个邻域 $W$ 和一个光滑函数 $\sigma:W\to\mathbb{R}^{n}$，使它在 $W\cap\varphi(U\cap V)$ 上与 $\rho^{-1}$ 一致。另一方面，因为我们假设 $\psi$ 是一个内图，存在一个以 $x_{0}$ 为中心、包含于 $\psi(U\cap V)$ 的开欧几里得球 $B$，所以 $\rho$ 本身在 $B$ 上是通常意义下光滑的。必要时缩小 $B$，我们可以假设 $B\subset\rho^{-1}(W)$。于是 $\sigma\circ\rho|_{B}=\rho^{-1}\circ\rho|_{B}=\operatorname{Id}_{B}$，所以由链式法则，对每个 $x\in B$，$D\sigma(\rho(x))\circ D\rho(x)$ 是恒等映射。由于 $D\rho(x)$ 是方阵，这蕴涵它非奇异。由推论~C.36 可得 $\rho$（作为从 $B$ 到 $\mathbb{R}^{n}$ 的映射）是一个开映射，所以 $\rho(B)$ 是 $\mathbb{R}^{n}$ 的一个开子集，它包含 $y_{0}=\varphi(p)$ 且包含于 $\varphi(V)$。这与 $\varphi(V)\subset H^{n}$ 且 $\varphi(p)\in\partial H^{n}$ 的假设矛盾。
\end{proof}

\end{document}